\subsection{5. Dieselbe Antwort auf demselben
Grundzustand}\label{dieselbe-antwort-auf-demselben-grundzustand}

Für Im z\textgreater0 und \(H_n=H|_{N=n}\):

\[\adjustbox{max width=\linewidth}{$\displaystyle 
G_{rs}(z)=\langle\Omega|f_r(z+E_0-H_{65})^{-1}f_s^\dagger|\Omega\rangle
+\langle\Omega|f_s^\dagger(z-E_0+H_{63})^{-1}f_r|\Omega\rangle.
$}\]

Die innere Symmetrie macht G diagonal und alle Diagonalelemente gleich.
Schreibe mit positiven Entnahme- und Additionsmaßen auf \(\epsilon>0\)

\[\adjustbox{max width=\linewidth}{$\displaystyle 
G(z)=\int\frac{d\nu_+(\epsilon)}{z-\epsilon}
+\int\frac{d\nu_-(\epsilon)}{z+\epsilon}.
$}\]

Dann gelten exakt

\[\adjustbox{max width=\linewidth}{$\displaystyle 
Z_+=\bar b/32,\quad Z_-=1-\bar b/32,
\quad a:=\int\epsilon\,d\nu_+=\int\epsilon\,d\nu_-
=\frac{\Delta\bar b-E_0}{64}.
$}\]

Für das gesamte signierte Spektralmaß:

\[\adjustbox{max width=\linewidth}{$\displaystyle 
\boxed{m_0=1,\quad m_1=0,\quad m_2=g^2S,\quad
m_3=g^2(\Delta S+7a).}
$}\]

Die ersten beiden Momente und die Kanalgewichte stimmen mit der
eingegangenen unabhängigen Rechnung überein. \textbf{Das dritte Moment
ist der zusätzliche eigene Schritt dieser Revision.}

\subsubsection{5.1 Herleitung des dritten
Moments}\label{herleitung-des-dritten-moments}

Setze \(\mathcal L A=[A,H]\). Stationarität macht diesen Operator
symmetrisch in der positiven, nach Nullvektoren quotientierten Metrik
\((A,B)=\langle\{A^\dagger,B\}\rangle\). Direkte Normalordnung ergibt

\[\adjustbox{max width=\linewidth}{$\displaystyle 
\sum_r\{[D_r,X],D_r^\dagger\}=-14Q_+.
$}\]

Die beiden Beiträge sind \(16Q_+\) und \(-30Q_+\). Ihre Koeffizienten
folgen aus den vollständig geprüften Kontraktionen
\(\sum_{rj}M_{A,rj}M_{B,rj}=16\delta_{AB}\) und
\(\sum_{Ar}M_{A,rj}M_{A,rk}=15\delta_{jk}\). Mit \([D,N_b]=D\) und
\(\langle Q_+\rangle=(E_0-\Delta\bar b)/(2g)\) folgt die Formel.

Als unabhängige Kontrolle wurden sämtliche Formeln einschließlich m3 an
der früher exakt geschlossenen N=4/N=5-Antwort symbolisch geprüft. Diese
Kontrolle verwendet den früheren Referenzzustand nur als Test der
Identitäten, nicht als Ersatz für den nativen Grundzustand.

Die allgemeine Methode, Spektralmomente aus Bewegungsgleichungen zu
gewinnen, ist etabliert; siehe
\href{https://arxiv.org/abs/0907.1284}{Freericks und Turkowski, Phys.
Rev.~B 80, 115119}. Die hier angegebenen W-Kontraktionen und Konstanten
wurden eigenständig für dieses Modell berechnet; die zitierte Arbeit
beweist keine TFPT-Aussage.

\subsubsection{5.2 Neue engere Dichte- und
Gewichtsgrenzen}\label{neue-engere-dichte--und-gewichtsgrenzen}

Positivität der beiden Antwort-Grammatrizen beziehungsweise zweimal
Cauchy-Schwarz liefern

\[\adjustbox{max width=\linewidth}{$\displaystyle 
m_2\ge a^2\left(\frac1{Z_-}+\frac1{Z_+}\right),
$}\] \[\adjustbox{max width=\linewidth}{$\displaystyle 
\boxed{(\Delta\bar b-E_0)^2\le
60g^2\bar b(32-\bar b)(1-7\bar b/480).}
$}\]

Mit der bereits bewiesenen Energieobergrenze und \(g/\Delta=1/20\) muss
das folgende rationale Polynom positiv sein:

\[\adjustbox{max width=\linewidth}{$\displaystyle 
P(b)=\frac7{3200}b^3-\frac{61}{50}b^2
+\frac{317591}{125000}b-\frac{79754843281}{62500000000}>0.
$}\]

Seine beiden im alten zulässigen Bereich liegenden Nullstellen liegen
bei ungefähr 0.842846697 und 1.245655664. Exakte rationale
Wurzelisolation ergibt die nach außen gerundete strenge Schranke

\[\adjustbox{max width=\linewidth}{$\displaystyle 
\boxed{0.842846<\bar b<1.245656.}
$}\]

Damit:

{\def\LTcaptype{none} % do not increment counter
\begin{longtable}[]{@{}
  >{\raggedright\arraybackslash}p{(\linewidth - 2\tabcolsep) * \real{0.4286}}
  >{\raggedleft\arraybackslash}p{(\linewidth - 2\tabcolsep) * \real{0.5714}}@{}}
\toprule\noalign{}
\begin{minipage}[b]{\linewidth}\raggedright
Größe pro Mode
\end{minipage} & \begin{minipage}[b]{\linewidth}\raggedleft
Strenges offenes Intervall
\end{minipage} \\
\midrule\noalign{}
\endhead
\bottomrule\noalign{}
\endlastfoot
Gesamtes Additionsgewicht & (0.0263389375, 0.03892675) \\
Gesamtes Entnahmegewicht & (0.96107325, 0.9736610625) \\
\(\langle\{\phi,\phi^\dagger\}\rangle\) & (0.981834183333\ldots,
0.987708495833\ldots) \\
Frühere \(\chi\)-Kompositnorm, nicht \(\phi\) & (0.040386370833\ldots,
0.059687683333\ldots) \\
\end{longtable}
}

Das sind Erwartungswerte und integrierte Spektralgewichte, keine
direkten Nachweise von Produktionsraten oder von bereits verfügbaren
Messinstrumenten.

\subsection{6. Der Polsatz und seine zusätzliche
Verschärfung}\label{der-polsatz-und-seine-zusuxe4tzliche-verschuxe4rfung}

Der eingegangene Beweis verwendet \(u_{k,r}=Q_+^kf_rF=f_rQ_+^kF\).
Invarianz und Besetzung ergeben

\[\adjustbox{max width=\linewidth}{$\displaystyle 
\langle u_{k,r},u_{k,s}\rangle=
\delta_{rs}\frac{64-2k}{64}\|Q_+^kF\|^2.
$}\]

Die Normen lauten 1, 465, 412200, 521164800, 833259571200. Die daraus
gebildete fünfdimensionale Variationsmatrix gibt 64 unabhängige
Richtungen unter \(-1.095812\Delta\). Das gesamte N=63-Komplement des
Nullbosonraums liegt über \(-3\Delta/4\). Minimax begrenzt den niedrigen
Raum auf genau 64 Dimensionen; seine injektive symmetrieverträgliche
Projektion auf die irreduzible duale 64 erzwingt ein einziges
Energieniveau.

Das beweist Existenz und Isolation. Seine Sichtbarkeit folgt aus den
positiven Entnahmemomenten. Mit

\[\adjustbox{max width=\linewidth}{$\displaystyle 
d=0.007737\Delta,\quad c=0.379636\Delta,
\quad a_{\max}=\frac{1.245656+1.158089}{64}\Delta
$}\]

gilt

\[\adjustbox{max width=\linewidth}{$\displaystyle 
Z_{\rm low}>\frac{c(1-1.245656/32)-a_{\max}}{c-d}
=\frac{40912436089}{46487375000}
=0.8800762806891118\ldots.
$}\]

Außerdem ist der niedrige Pol die kleinste Entnahmeenergie. Daher
\(a\ge\epsilon_{\rm low}Z_-\), also bereits ohne vollständige
Polauswertung

\[\adjustbox{max width=\linewidth}{$\displaystyle 
\frac{\epsilon_{\rm low}}\Delta
<\frac{1.245656+1.158089}{64-2(1.245656)}
=\frac{2403745}{61508688}=0.039079763821332\ldots.
$}\]

Diese Verschärfungen entstehen aus dem \textbf{Zusammenführen
kompatibler Beweise}, nicht aus einem neuen gewählten Parameter. Das
Restgewicht des gesamten CAR-Maßes ist somit kleiner als
0.119923719311\ldots{} .

\subsection{7. Einfache Organisation, aber keine falsche
Zwei-Linien-Lösung}\label{einfache-organisation-aber-keine-falsche-zwei-linien-luxf6sung}

Die ersten zwei orthonormalen Operatoren sind f und \(D/\sqrt S\). Die
erste Resolventenreduktion hat deshalb die exakte Form

\[\adjustbox{max width=\linewidth}{$\displaystyle 
\boxed{G(z)=\frac1{z-\displaystyle\frac{g^2S}{z-a_1-\Sigma_2(z)}}},
\qquad a_1=\Delta+\frac{7a}{S}.
$}\]

\(\Sigma_2\) ist die Resolvente des verbleibenden nativen Operatorraums,
gekoppelt an den dazu orthogonalen Rest von \(\mathcal L(D/\sqrt S)\).
Es wurde kein äußeres Bad hinzugefügt. Dies ist eine genaue Ordnung der
Rechnung, \textbf{keine abgeschlossene Berechnung von \(\Sigma_2\)} und
keine Vereinigung sämtlicher TOE-Aufgaben in einer einzigen bereits
gelösten Funktion.

\subsubsection{7.1 Warum man die Rückwirkung nicht exakt weglassen
darf}\label{warum-man-die-ruxfcckwirkung-nicht-exakt-weglassen-darf}

Angenommen, es gäbe nur je eine Entnahme- und Additionslinie mit
Energien \(\epsilon_-,\epsilon_+\). Symmetrie macht diese für alle r
gleich. Dann liefern die Bewegungsgleichungen auf demselben Grundzustand

\[\adjustbox{max width=\linewidth}{$\displaystyle 
D_r\Omega=-\frac{\epsilon_-}{g}f_r\Omega,
\qquad D_r^\dagger\Omega=\frac{\epsilon_+}{g}f_r^\dagger\Omega.
$}\]

Summe nach Multiplikation mit \(f_r^\dagger\) beziehungsweise \(f_r\)
und N=64 ergeben

\[\adjustbox{max width=\linewidth}{$\displaystyle 
H\Omega=\left[-32\epsilon_-+
(\Delta-\epsilon_++\epsilon_-)N_b\right]\Omega.
$}\]

Ein Eigenzustand mit negativer Energie kann keine feste Bosonenzahl
haben: Dann wäre \(\langle Q_++Q_-\rangle=0\) und seine Energie
\(\Delta\langle N_b\rangle\ge0\). Daher sind \(\Omega\) und
\(N_b\Omega\) unabhängig. Die angenommene Zweilinienform erzwingt
\(\epsilon_+-\epsilon_-=\Delta\).

Andererseits geben \(m_0=1,m_1=0\) für ein Zweilinienmaß
\(m_3/m_2=\epsilon_+-\epsilon_-\). Die exakt bestimmte Formel lautet
jedoch

\[\adjustbox{max width=\linewidth}{$\displaystyle 
\frac{m_3}{m_2}=\Delta+7a/S>\Delta.
$}\]

Widerspruch. Damit ist \(\Sigma_2\not\equiv0\) bewiesen. Mindestens ein
Kanal besitzt mehr als eine Energielinie. Der dominante isolierte
Entnahmepol und diese unvermeidliche Reststruktur widersprechen einander
nicht. Eine Zweilinienform könnte höchstens eine zu zertifizierende
Näherung sein.

\subsection{8. Was der Operationssatz jetzt tatsächlich
hergibt}\label{was-der-operationssatz-jetzt-tatsuxe4chlich-hergibt}

{\def\LTcaptype{none} % do not increment counter
\begin{longtable}[]{@{}
  >{\raggedright\arraybackslash}p{(\linewidth - 4\tabcolsep) * \real{0.3333}}
  >{\raggedright\arraybackslash}p{(\linewidth - 4\tabcolsep) * \real{0.3333}}
  >{\raggedright\arraybackslash}p{(\linewidth - 4\tabcolsep) * \real{0.3333}}@{}}
\toprule\noalign{}
\begin{minipage}[b]{\linewidth}\raggedright
Vertrag
\end{minipage} & \begin{minipage}[b]{\linewidth}\raggedright
Mathematisch abgesichert
\end{minipage} & \begin{minipage}[b]{\linewidth}\raggedright
Nicht dadurch verfügbar
\end{minipage} \\
\midrule\noalign{}
\endhead
\bottomrule\noalign{}
\endlastfoot
Markierter endlicher Matrixcompiler & Matrizen, Ordnungen, konkrete
endliche Syntheseidentitäten & Vollständiges Fockinstrument,
Präparation, physischer Zeitgenerator \\
H allein & Modellzeitentwicklung & Unabhängiges Schalten von X und
\(N_b\) \\
X und \(N_b\), wenn als Kontrollen gewährt & N=3-Kontrollalgebra der
Dimension 14 & Beliebige Modenoperationen \\
Zusätzlich dokumentierter vorzeichenrichtiger Clock-Lift &
N=3-Kontrollalgebra der Dimension \textbf{84} & Vollständige
Zustandsunterscheidung oder Ladung-eins-Instrument \\
\end{longtable}
}

Die neue Clock-Rechnung wurde übernommen und reproduziert. Die fünf
Multiplizitätszeilen für die sechs Clockphasen sind:

{\def\LTcaptype{none} % do not increment counter
\begin{longtable}[]{@{}lrrrrrr@{}}
\toprule\noalign{}
Teilraum & 0 & 1 & 2 & 3 & 4 & 5 \\
\midrule\noalign{}
\endhead
\bottomrule\noalign{}
\endlastfoot
Dunkle F3 & 6384 & 6280 & 6280 & 6384 & 6280 & 6280 \\
Dunkle BF & 16 & 8 & 8 & 16 & 8 & 8 \\
Aktiv 7 & 560 & 440 & 440 & 560 & 440 & 440 \\
Aktiv 10 & 112 & 88 & 88 & 112 & 88 & 88 \\
Aktiv 12 & 80 & 40 & 40 & 80 & 40 & 40 \\
\end{longtable}
}

Die aktiven Zeilen tragen zusätzlich einen Zweiniveaufaktor. Alle 45504
N=3-Zustände sind erfasst; der Kommutant hat weiterhin Dimension
240742144. Die alte Zahl 14 beschreibt den engeren Zweikontrollvertrag,
nicht den jetzt zusätzlich geprüften Clockvertrag.

\subsubsection{8.1 Konkrete
Präparationsgrenze}\label{konkrete-pruxe4parationsgrenze}

Alle genannten Fockkontrollen erhalten N. Ein Wort aus ihnen und jeder
einzeln zahlenerhaltende ausgewählte Messzweig bleiben im
Ausgangssektor. Sie können aus dem leeren N=0-Zustand \textbf{nicht} den
N=64-Grundzustand erzeugen. Diese Aussage betrifft zahlenerhaltende
Krauszweige; eine bloß U(1)-kovariante offene Dynamik kann sehr wohl
geladene Krausoperatoren haben.

Auch aus dem voll besetzten F führt bloßes \(e^{-itH}\) nicht zur
Konvergenz auf \(\Omega\): Der Grundzustandsüberlapp bleibt konstant.
Normiertes \(e^{-\tau H}F\) konvergiert mathematisch mit der bewiesenen
Lücke, benötigt aber eine gesonderte Instrument- und
Ressourcenherleitung.

\subsubsection{8.2 Der kleinste konkrete
Ladungstest}\label{der-kleinste-konkrete-ladungstest}

Ein Kandidat ist eine explizite Referenzmode c mit
\(T=\lambda(c^\dagger f_r+f_r^\dagger c)\). Sie erhält die gemeinsame
Ladung, verändert aber die Ladung der ursprünglichen Bank. Der endliche
Austausch- und Detuningtest ist exakt; \textbf{dieses T wurde nicht als
neuer nativer Hamiltonterm eingeführt}.

Für dieselbe N=64-Präparation unter \(H+\mu N\) gilt
\(G_\mu(z)=G_0(z-\mu)\), insbesondere \(m_1=\mu\). Das misst nur gegen
eine festgelegte Referenz: Ein gemeinsamer Zusatz
\(\mu(N_{\rm Bank}+N_{\rm Referenz})\) bleibt unsichtbar. Die Spektral-
schranken dieses Berichts gelten als Grundzustandsschranken nur für μ=0.

\subsection{9. Relativistisches Wörterbuch: Ausschluss und zwei präzise
Alternativen}\label{relativistisches-wuxf6rterbuch-ausschluss-und-zwei-pruxe4zise-alternativen}

Die Indexfrage darf nicht nachträglich durch einen Namen beantwortet
werden. Die Konventionen für Zweikomponentenspinoren wurden mit
\href{https://arxiv.org/abs/0812.1594}{Dreiner, Haber und Martin}
abgeglichen; die folgenden speziellen W-Tests sind eigene exakte
Tensorrechnungen.

\subsubsection{9.1 Was nicht funktioniert}\label{was-nicht-funktioniert}

Für 64 gleichhändige Weylfelder mit innerem Index I und einen skalaren
Vermittler ist \(M_A\otimes\varepsilon_{\rm Lorentz}\) symmetrisch. Die
Grassmann-Antisymmetrisierung verschwindet daher identisch, für alle 60
W-Zeilen. Außerdem lässt die volle irreduzible innere Darstellung
\(16\otimes4\) auf denselben 64 Komponenten nach Schur keine zusätzliche
kommutierende nichttriviale Lorentzspinorwirkung zu.

\subsubsection{9.2 Bereits bekannter nichtverschwindender
Typ}\label{bereits-bekannter-nichtverschwindender-typ}

Ein symmetrischer Lorentzspinortensor koppelt an das antisymmetrische M.
Der Vermittler hat dann einen passenden dualen (1,0)-Typ samt
adjungiertem Typ. Diese Variante besteht den algebraischen
Nichtnulltest, aber noch keinen vollständigen Kinetik-, Positivitäts-,
Constraints- oder Herkunftstest.

\subsubsection{9.3 Neue minimale skalare Alternative - ausdrücklich
zusätzlich}\label{neue-minimale-skalare-alternative---ausdruxfccklich-zusuxe4tzlich}

Will man zugleich W, gleiche Weylhand, lokale Bilinearität und einen
skalaren Vermittler behalten, kann man einen \textbf{unabhängigen}
Hilfsindex a=1,2 mit alternierender Form einführen:

\[\adjustbox{max width=\linewidth}{$\displaystyle 
b_A^\dagger(M_A)_{IJ}\varepsilon_{ab}
\varepsilon_{\alpha\beta}\psi_{I a\alpha}\psi_{J b\beta}
+\mathrm{h.c.}
$}\]

M und \(\varepsilon_{ab}\) sind jeweils antisymmetrisch; ihr Produkt ist
als innerer Kopplungstensor symmetrisch. Zusammen mit der Lorentz-
Epsilonform ist der Gesamttensor antisymmetrisch und nicht null. Alle 60
Kanäle bestehen den Test. Eine nichtverschwindende alternierende Form
existiert nicht in Dimension eins; in Dimension zwei ist sie bis auf
Normierung eindeutig. \textbf{In dieser eng benannten Klasse von
Tensorfaktor-Reparaturen ist die binäre Ergänzung minimal.}

Das ist noch keine gefundene native Lösung: Sie verdoppelt die inneren
Weylkomponenten von 64 auf 128, zusätzlich zu deren Lorentzspinorindex.
Ein Double-Cover-Minuszeichen stellt nicht automatisch zwei unabhängige
Felder bereit. Auch Nambu-Umbenennung \((f,f^\dagger)\) genügt nicht:
Das gemischte Produkt hat Ladung null statt −2, sodass derselbe
\(b^\dagger ff\)-Ladungsvertrag nicht erhalten bleibt.

Der entscheidende Quellenauftrag ist daher eng: Gibt es diese zweite
gleichgeladene, CAR-unabhängige Komponente bereits im tatsächlichen
Compilerprozess? Und liefert ihre Projektion genau das bisherige H und
die geprüfte Antwort? Ohne beides bleibt die skalare Alternative ein
zusätzliches Modell, nicht eine Erklärung des ursprünglichen.

